\diarytitle{0083}{4点データに対する直線の最小二乗フィット}

データ点 $(x_i, y_i)$ ($i=1,\dots,n$) に対して、モデル関数 $f(x)$（未知パラメータを含む）を当てはめる最小二乗法では、誤差二乗和
\[
S = \sum_{i=1}^{n} \bigl(y_i - f(x_i)\bigr)^2
\]
を最小にするようにパラメータを定める。$f$ が未知パラメータについて線形（ここでは $f(x) = ax + b$）である場合、$S$ は各パラメータの2次関数であるから、最小値では各パラメータに関する偏微分がすべて $0$ になる。

\[
S(a,b) = \sum_{i=1}^n (y_i - ax_i - b)^2
\]
とおくと、
\[
\frac{\partial S}{\partial a} = -2\sum_i x_i(y_i - ax_i - b) = 0, \qquad
\frac{\partial S}{\partial b} = -2\sum_i (y_i - ax_i - b) = 0
\]
これを整理すると、正規方程式
\[
\left\{
\begin{aligned}
a\sum_i x_i^2 + b \sum_i x_i &= \sum_i x_i y_i \\
a\sum_i x_i + b\, n &= \sum_i y_i
\end{aligned}
\right.
\]
を得る。

$n=4$ とし、必要な和を計算すると
\[
\sum x_i = 6,\quad \sum x_i^2 = 14,\quad \sum y_i = 9,\quad \sum x_i y_i = 18
\]
正規方程式は
\[
14a + 6b = 18, \qquad 6a + 4b = 9
\]
第1式を $2$ で割って $7a+3b=9$。これを $4$倍、第2式を $3$倍すると
\[
28a+12b=36, \qquad 18a+12b=27
\]
辺々引いて $10a = 9$、すなわち $a = \dfrac{9}{10}$。これを $6a+4b=9$ に代入すると
\[
6\cdot\frac{9}{10} + 4b = 9 \ \Longrightarrow\ 4b = 9 - \frac{27}{5} = \frac{18}{5} \ \Longrightarrow\ b = \frac{9}{10}
\]
したがって求める直線は
\[
\boxed{\; y = \frac{9}{10}x + \frac{9}{10} \;}
\]

（検算: 残差 $r_i = y_i - (0.9x_i+0.9)$ は $r_1=0.1,\ r_2=0.2,\ r_3=-0.7,\ r_4=0.4$ で、$\sum r_i = 0$ かつ $\sum x_i r_i = 0\cdot0.1+1\cdot0.2+2\cdot(-0.7)+3\cdot0.4=0.2-1.4+1.2=0$ となり、正規方程式を満たすことが確認できる。）
